\section{Baseline cuối của hệ thống}

Phần này chỉ chứa nội dung có phạm vi toàn hệ thống. Bản cuối được lấy từ \path{docs/phase-1/group-report}, cập nhật các thay đổi hợp lệ phát sinh ở Phase 2--3 và loại bỏ nội dung không còn đúng với MVP.

\subsection{Tổng quan và đặc tả dự án}
Giữ mối liên hệ từ bối cảnh và vấn đề đến stakeholder need, user story, mục tiêu, phạm vi và system boundary.

\subsection{Yêu cầu hệ thống}
Giữ mã requirement ổn định; đánh dấu rõ requirement đã hiện thực, chưa hiện thực hoặc thay đổi có phê duyệt. Không sửa requirement chỉ để khớp với code mà không giải thích.

\subsection{Actor Catalog và Use-case Model toàn hệ thống}
Giữ actor, external system và use-case diagram toàn hệ thống. Đây là bản đồ định hướng trước khi người đọc đi vào bảy phân hệ; đặc tả chi tiết của 18 use case được đặt cùng artefact thiết kế tương ứng ở phần sau.

\subsection{Quy tắc tích hợp và chống lặp}

Context diagram, use-case diagram toàn hệ thống, deployment diagram và component diagram chỉ xuất hiện một lần. Mockup, class diagram hoặc state chart dùng chung được đặt tại phân hệ/UC sở hữu chính; các UC khác dẫn chiếu bằng mã và số hình. Mọi artefact phải giữ ổn định các mã \texttt{SUB-*}, \texttt{UC-*}, \texttt{FR-*}, \texttt{BR-*}, \texttt{SCR-*}, \texttt{SEQ-*}, \texttt{ACT-*}, \texttt{STM-*}, \texttt{CD-*} và \texttt{TC-*}.
